\chapter{HbA1c Test}
\index{HbA1c Test}

\section*{Definition and Overview}
The HbA1c test, also called glycated haemoglobin, is a blood test that measures the average blood glucose level over the previous two to three months. Haemoglobin is the oxygen-carrying protein in red blood cells, and glucose attaches to it in proportion to blood glucose levels; the more glucose in the blood, the more "glycated" the haemoglobin becomes. Because red blood cells live for about three months, the HbA1c reflects long-term glucose control rather than a single day's reading. It is the key test for diagnosing and monitoring diabetes, and it guides treatment decisions. HbA1c is reported as a percentage or in mmol/mol.

\section*{Epidemiology}
Diabetes affects around 1 in 10 people, and the HbA1c test is central to its diagnosis and management. Around 5--10% of people have HbA1c levels in the prediabetes range, indicating an increased risk of developing diabetes. The test is also used to monitor the effectiveness of treatment in the millions of people with type 1 and type 2 diabetes. Because HbA1c does not require fasting, it is convenient and widely used in both laboratory and point-of-care settings.

\section*{Aetiology and Pathophysiology}
The HbA1c test measures the average of glucose exposure over the lifespan of red blood cells.

\begin{itemize}
    \item \textbf{Glycation:} glucose in the blood binds non-enzymatically to haemoglobin A, forming glycated haemoglobin (HbA1c)
    \item \textbf{Reflects average glucose:} the level of HbA1c is proportional to the average blood glucose over the previous 8--12 weeks
    \item \textbf{Red cell lifespan:} because red blood cells survive around 120 days, the test reflects weeks of glucose exposure, not just the day of the test
    \item \textbf{Diagnosis:} an HbA1c of 6.5% (48 mmol/mol) or higher diagnoses diabetes; 6.0--6.4% indicates prediabetes in local guidelines
    \item \textbf{Monitoring:} targets for people with diabetes are typically below 7% (53 mmol/mol), individualised to the person
    \item \textbf{Limitations:} conditions affecting red cell lifespan or haemoglobin, such as anaemia and recent transfusion, can affect results
\end{itemize}

\section*{Clinical Features}
The HbA1c test is used in specific clinical situations.

\begin{itemize}
    \item Diagnosing type 2 diabetes and prediabetes
    \item Monitoring glucose control in people with known diabetes
    \item Assessing the effectiveness of diet, exercise, and medicines
    \item Deciding when to start or adjust diabetes medicines
    \item Screening people at high risk of diabetes
    \item Assessing control in pregnancy-related diabetes in some settings (though other tests are preferred)
    \item Guiding treatment targets in individualised care
\end{itemize}

\section*{Differential Diagnosis}
Results should be interpreted with awareness of factors that affect HbA1c.

\begin{itemize}
    \item Anaemia and haemoglobinopathies, which can lower or raise HbA1c inaccurately
    \item Recent blood transfusion
    \item Kidney failure and pregnancy, which affect interpretation
    \item Rapid changes in glucose, which may not yet be reflected in HbA1c
    \item Conditions that shorten red cell survival
    \item Laboratory variation between methods
    \item Fasting glucose and oral glucose tolerance tests, which may be needed when HbA1c is not reliable
\end{itemize}

\section*{When to Seek Medical Care}
People at risk of diabetes should be screened regularly, and people with diabetes should have HbA1c measured at least every three to six months. See a doctor promptly for symptoms of high blood glucose, including increased thirst, frequent urination, fatigue, and blurred vision.

\textbf{Contact your local emergency services for:}
\begin{itemize}
    \item This test is not an emergency; however, seek urgent care for very high blood glucose with vomiting, confusion, or drowsiness
    \item Signs of diabetic ketoacidosis, including rapid breathing and fruity breath
    \item Very low blood glucose with confusion or collapse
    \item Any medical emergency related to diabetes
\end{itemize}

\section*{Diagnosis and Investigations}
The HbA1c test is performed on a small blood sample.

\begin{itemize}
    \item \textbf{Blood sample:} a small amount of blood from a vein or finger prick
    \item \textbf{No fasting required:} the test can be done at any time of day
    \item \textbf{Repeat testing:} an elevated result is confirmed with a second test unless symptoms are present
    \item \textbf{Monitoring schedule:} every three to six months for people with diabetes
    \item \textbf{Laboratory or point-of-care testing:} both are available
    \item \textbf{Interpretation:} results are interpreted by a doctor in the context of the whole picture
\end{itemize}

\section*{Management}
The HbA1c result guides diabetes management.

\begin{itemize}
    \item \textbf{Setting targets:} a personalised target, often below 7%, is agreed between the person and their doctor
    \item \textbf{Lifestyle measures:} diet, physical activity, and weight management lower HbA1c
    \item \textbf{Medicines:} metformin and other medicines are started and adjusted based on HbA1c
    \item \textbf{Insulin:} intensification or adjustment of insulin therapy for type 1 and advanced type 2 diabetes
    \item \textbf{Monitoring glucose:} home blood glucose monitoring complements HbA1c
    \item \textbf{Reviewing complications:} regular checks for eye, kidney, and foot complications in people with diabetes
    \item \textbf{Education:} diabetes education supports self-management
    \item \textbf{Regular retesting:} ongoing measurement tracks control over time
\end{itemize}

\section*{Complications}
\begin{itemize}
    \item Misinterpretation of results when conditions affecting red cells are not considered
    \item A false sense of security from a single HbA1c that misses glucose variability
    \item Under-treatment if targets are too high, or over-treatment with hypoglycaemia risk if too low
    \item Delayed diagnosis of diabetes when screening is not done
    \item Complications of diabetes if control remains poor
    \item Cost and inconvenience of repeated testing
\end{itemize}

\section*{Prognosis and Outcomes}
The HbA1c test has transformed diabetes care. Every 1% reduction in HbA1c is associated with a significant reduction in the risk of diabetes complications, including eye, kidney, and nerve disease and heart attack. People who achieve their HbA1c targets substantially reduce their risk of complications and maintain better health and quality of life. The test also enables earlier diagnosis of prediabetes and diabetes, allowing intervention before complications develop.

\section*{Prevention and Self-Care}
\begin{itemize}
    \item Have regular HbA1c tests if you have diabetes or risk factors for it
    \item Know your target and your current level
    \item Follow your treatment plan and attend appointments
    \item Maintain a healthy diet, exercise regularly, and manage weight
    \item Take medicines as prescribed
    \item Monitor blood glucose at home as advised
    \item Report new or worsening symptoms promptly
    \item If you have anaemia or other conditions affecting red cells, discuss how this affects your test
\end{itemize}

\section*{Sources and Further Reading}
The following reputable sources provide further information for healthcare students and patients seeking reliable, current advice about the HbA1c test.

\begin{itemize}
    \item Diabetes Australia. \textit{HbA1c: understanding your result.} https://www.diabetesaustralia.com.au/
    \item Healthdirect Australia. \textit{HbA1c test.} https://www.healthdirect.gov.au/
    \item Royal College of Pathologists of Australasia. \textit{Patient information on HbA1c.} https://www.rcpa.edu.au/
    \item Australian Government Department of Health and Aged Care. \textit{Diabetes information.} https://www.health.gov.au/
    \item NHS. \textit{HbA1c test for diabetes.} https://www.nhs.uk/
    \item American Diabetes Association. \textit{Understanding A1C.} https://diabetes.org/
\end{itemize}
